% ===========================================================================
% reviewmarks.tex -- relevance colour-coding and margin commentary.
%
% TEMPORARY REVIEW APPARATUS. Lives on branch review-colour-map only; it is
% not part of the submitted manuscript and must not be merged into a
% submission branch without first setting \reviewmarksfalse (or removing the
% \input from main.tex and the \rb/\re markers from the sections).
%
% ---------------------------------------------------------------------------
% WHAT IT DOES
%
% Every prose paragraph and every figure/table caption in the body and the
% appendices is wrapped as
%
%     \rb{<class>}{<ID>}{<two-sentence rationale>}%
%     ... the paragraph, untouched ...\re
%
% \rb opens a group, colours everything up to the matching \re according to
% <class>, prints the ID inline as a small coloured tag, and emits a margin
% note carrying the ID, the class and the rationale. \re closes the group.
% The manuscript text itself is never edited by this apparatus: turning the
% master switch off restores the manuscript exactly.
%
% Five classes, per the request that drove this branch:
%
%   K  green   keep      must stay in the body as it is
%   T  orange  trim      could be removed; the paper survives without it
%   A  purple  annex     belongs in an appendix, not the body
%   S  brown   supp      belongs in the supplementary material
%   D  red     delete    can be deleted outright, nothing is lost
%
% ---------------------------------------------------------------------------
% SWITCHES
%
%   \reviewmarksfalse    MASTER. Off = clean manuscript, no colour, no notes,
%                        no tags, normal JFM page size. All other switches
%                        are then irrelevant.
%   \reviewcolourfalse   keep the notes and tags, drop the text colouring
%   \reviewnotesfalse    keep the colouring, drop the margin notes
%   \reviewidsfalse      keep everything, drop the inline [ID] tags
%
% Flip them in the block marked SWITCHES below.
%
% ---------------------------------------------------------------------------
% PAGE SIZE
%
% The JFM trim (174mm) leaves 21.5mm of right margin, which cannot hold a
% legible note. In review mode the paper grows by \reviewextrawidth and all of
% the extra goes to the right margin; the text block keeps its exact position
% and its measure, so line breaking is unaffected by the widening itself.
% \@twosidefalse and the equalised side margins force every note into the right
% margin (marginnote picks its side from \if@twoside).
%
% Nothing below is JFM-specific: the widening is relative to whatever page size
% the class set. To retune, \providecommand any of these BEFORE \input:
%
%   \reviewextrawidth   how much wider the page becomes   (default 62mm)
%   \reviewnotewidth    margin-note measure               (default 46mm)
%   \reviewnotegutter   gap between text block and note   (default 5mm)
% ===========================================================================

\usepackage{marginnote}

\providecommand{\reviewextrawidth}{62mm}
\providecommand{\reviewnotewidth}{46mm}
\providecommand{\reviewnotegutter}{5mm}

% This file is \input from the main.tex preamble, where @ is not a letter, so
% the internal \rv@... macros need the catcode change for the whole file.
\makeatletter

% ------------------------------- SWITCHES ---------------------------------
\newif\ifreviewmarks   \reviewmarkstrue    % MASTER
\newif\ifreviewcolour  \reviewcolourtrue
\newif\ifreviewnotes   \reviewnotestrue
\newif\ifreviewids     \reviewidstrue
% --------------------------------------------------------------------------

% Colours are defined unconditionally (cheap, and \pending already uses red).
% The class preloads `color`, so \definecolor is available without xcolor.
\definecolor{rvK}{rgb}{0.00,0.42,0.16}   % green   keep
\definecolor{rvT}{rgb}{0.78,0.42,0.00}   % orange  trim
\definecolor{rvA}{rgb}{0.44,0.12,0.60}   % purple  annex
\definecolor{rvS}{rgb}{0.42,0.26,0.09}   % brown   supplementary
\definecolor{rvD}{rgb}{0.78,0.06,0.06}   % red     delete
\definecolor{rvGrey}{rgb}{0.35,0.35,0.35}

\expandafter\def\csname rvname@K\endcsname{KEEP}
\expandafter\def\csname rvname@T\endcsname{TRIM?}
\expandafter\def\csname rvname@A\endcsname{ANNEX}
\expandafter\def\csname rvname@S\endcsname{SUPPL.}
\expandafter\def\csname rvname@D\endcsname{DELETE}

% Deferred to \AtBeginDocument: lineno-FLM.sty, which the JFM class loads for
% the `lineno' option, ends by forcing \marginparwidth to 2.6pc, and it is
% loaded after this file is read. Setting the lengths here instead is what
% makes the notes 46mm wide rather than 11mm.
\ifreviewmarks
  \AtBeginDocument{%
    \@twosidefalse                        % all notes to the right margin
    \setlength\evensidemargin{\oddsidemargin}%
    \setlength\marginparwidth{\reviewnotewidth}%
    \setlength\marginparsep{\reviewnotegutter}%
    \setlength\marginparpush{4\p@}%
    \renewcommand*{\marginfont}{\normalfont}%
  }
  \pdfpagewidth=\dimexpr\pdfpagewidth+\reviewextrawidth\relax
  \special{papersize=\the\pdfpagewidth,\the\pdfpageheight}
\fi

% The margin note: a coloured rule, the ID and class, then the rationale.
\newcommand{\rv@note}[3]{%
  \marginnote{%
    \raggedright\setlength{\parindent}{0pt}\normalfont
    {\color{rv#1}\rule{\marginparwidth}{1.1pt}}\par\nobreak\vspace{1.5pt}
    {\scriptsize\bfseries\color{rv#1}#2\hfill\csname rvname@#1\endcsname}\par
    \nobreak\vspace{1pt}
    {\scriptsize\color{rvGrey}#3\par}%
  }%
}

% The inline tag, so an ID is still readable when the notes are switched off
% and so a colour can be traced to a row of the ledger from the text itself.
% It inherits the class colour from \rb, which sets it first.
\newcommand{\rv@tag}[2]{{\scriptsize [#2]}\,}

% \rb{class}{id}{note} ... \re
% The \fi sits inside the group deliberately: TeX does not require \begingroup
% and \endgroup to nest with conditionals, and this keeps both macros
% single-token no-ops when the master switch is off.
% Both are robust: the caption markers sit inside \caption{}, which the caption
% package double-scans and writes to main.lof, and a fragile \if... there breaks
% brace balance ("Argument of \caption@ydblarg has an extra }").
% Two details here are load-bearing, both found by bisection after the captions
% came out coloured and the body prose stayed black:
%   \leavevmode  -- at the head of a body paragraph \rb runs in vertical mode,
%                   where \color does not reach the paragraph that follows.
%                   Captions worked from the start because they are already in
%                   horizontal mode.
%   \color twice -- \marginnote typesets its note in a box whose own colour
%                   push and pop leave the surrounding text back at black, so
%                   the class colour has to be re-asserted after the note. The
%                   two pushes pair with two \reset@color at \re's \endgroup.
\DeclareRobustCommand{\rb}[3]{%
  \ifreviewmarks
    \begingroup
    \leavevmode
    \ifreviewcolour\color{rv#1}\fi
    \ifreviewnotes\rv@note{#1}{#2}{#3}\fi
    \ifreviewcolour\color{rv#1}\fi
    \ifreviewids\rv@tag{#1}{#2}\fi
  \fi
  \ignorespaces
}
\DeclareRobustCommand{\re}{%
  \ifreviewmarks
    \endgroup
  \fi
}

% Legend, placed once at the head of the body.
\newcommand{\rv@key}[2]{%
  \raisebox{0.5ex}{\color{rv#1}\rule{9pt}{4pt}}~{\color{rv#1}\textbf{#1}}~#2%
}
\newcommand{\reviewlegend}{%
  \ifreviewmarks
    \begingroup
    \small\noindent
    \fbox{\parbox{\dimexpr\textwidth-2\fboxsep-2\fboxrule\relax}{%
      \footnotesize\setlength{\parindent}{0pt}%
      \textbf{Editorial overlay (not part of the manuscript).} Every paragraph
      and caption is colour-coded for its value to the paper and carries a
      unique identifier; the margin note gives the reason. Switch the overlay
      off with \texttt{\textbackslash reviewmarksfalse} in
      \texttt{reviewmarks.tex}.\par\vspace{3pt}
      \rv@key{K}{keep in the body.}\quad
      \rv@key{T}{could be removed.}\par\vspace{1pt}
      \rv@key{A}{move to an appendix.}\quad
      \rv@key{S}{move to the supplementary material.}\quad
      \rv@key{D}{delete.}\par\vspace{3pt}
      Identifiers are \texttt{<unit>-<nn>}: \texttt{AB} abstract, \texttt{S1}
      to \texttt{S6} the numbered sections (\texttt{S31}, \texttt{S34} and so
      on for the \texttt{v4} subfiles), \texttt{APA} to \texttt{APC} the
      appendices; \texttt{C} in the number marks a caption. The full ledger,
      with the same identifiers, is
      \texttt{editorial/RELEVANCE\_MAP.md}.%
    }}\par\vspace{6pt}
    \endgroup
  \fi
}

\makeatother
